\documentclass[11pt,a4paper]{article}
\usepackage[hyperref]{acl2023}
\usepackage{times}
\usepackage{latexsym}
\usepackage{amsmath,amssymb,amsthm}
\usepackage{booktabs}
\usepackage{graphicx}
\usepackage{multirow}
\usepackage{xcolor}
\usepackage{algorithm}
\usepackage{algpseudocode}
\usepackage{subcaption}

\definecolor{privred}{RGB}{231,76,60}
\definecolor{utilblue}{RGB}{52,152,219}
\definecolor{novelgreen}{RGB}{46,204,113}

\title{Privacy-Preserving Text Anonymization via Decoupled Mask-and-Fill:\\Adversarial Probing, Adaptive Privacy Budgets, and Cross-Lingual Transfer}

\author{
  Anonymous Submission
}

\begin{document}
\maketitle

% ============================================================================
\begin{abstract}
We propose a \emph{decoupled mask-and-fill} architecture for
privacy-preserving text anonymization that separates PII detection from
replacement generation across disjoint data partitions.  The AI4Privacy~500K
multilingual dataset is split into two language-stratified halves: Half-A
trains a \textbf{Censor} (NER token-classifier) and Half-B trains a
\textbf{Hallucinator} (seq2seq generator), ensuring the generator never
observes real entities.  Our baseline employs DeBERTa-v3-base and
Flan-T5-Large with QLoRA; the advanced variant upgrades to XLM-RoBERTa
and mT5-large with DP-SGD and entity-consistency hashing.
We identify a critical limitation of entity-replacement approaches:
\emph{contextual re-identification}—even after swapping entities, the
surrounding context (superlatives, unique roles, specific amounts) can
re-identify the individual.  To address this, we propose two additional models:
(3)~a \textbf{Context-Aware Rephraser} that adds a third stage to the decoupled
pipeline, taking entity-replaced text and generalising contextual signals; and
(4)~an \textbf{End-to-End Semantic Privacy Paraphraser} that performs entity
replacement and contextual generalisation in a single pass.
Building on this four-model foundation, we introduce three novel techniques:
(1)~\textbf{Adversarial Privacy Probing (APP)}, which uses gradient reversal to
force the encoder to produce representations from which original entities
cannot be reconstructed;
(2)~\textbf{Entity-Type Adaptive Privacy (ETAP)}, the first heterogeneous
per-entity-type differential privacy budget allocation for text anonymization; and
(3)~\textbf{Cross-Lingual Privacy Transfer (CLPT)}, language-adversarial
training that enables zero-shot privacy-preserving anonymization in unseen
languages.
For comparison, we also evaluate five end-to-end seq2seq baselines
(T5-Efficient-Tiny through DistilBART), finding persistent entity
leakage (1.6--3.0\%) and entity-consistency failures.
APP reduces entity reconstruction accuracy by 72\%, ETAP improves
low-sensitivity entity utility by 8--12~F1 points without degrading
critical-entity protection, CLPT achieves 89\% of supervised performance
in zero-shot cross-lingual transfer, and the Semantic Privacy Paraphraser
reduces contextual re-identification risk by over 60\% compared to
entity-replacement alone.
\end{abstract}

% ============================================================================
\section{Introduction}
\label{sec:intro}

Text anonymization is a prerequisite for deploying NLP models on sensitive data
in healthcare, finance, and legal domains.  The dominant paradigm follows either
a \emph{detect-then-replace} pipeline—where an NER model identifies PII and a
generator produces replacements \cite{lison2021anonymization,
pilán2022textanonymization}—or an \emph{end-to-end seq2seq} approach that
directly maps original text to anonymized versions.

We argue that the \emph{decoupled mask-and-fill} paradigm offers a
fundamental architectural advantage: by training the NER censor and the
replacement generator on \emph{disjoint} data partitions, the generator
never observes real entities, providing privacy by construction.  We
partition the AI4Privacy~500K multilingual dataset into two
language-stratified halves and train a Censor (NER) on Half-A and a
Hallucinator (seq2seq) on Half-B.

For comparison, end-to-end seq2seq models (T5-Efficient-Tiny through
DistilBART-CNN) achieve BLEU up to 69.53 but exhibit persistent entity
leakage (1.6--3.0\%), entity-consistency breakdowns, and
size-dependent pathologies (DistilBART-CNN catastrophically degenerates at
333M parameters).

These findings—both from the seq2seq baselines and from the decoupled
pipeline—expose three fundamental limitations that current approaches
fail to address:

\paragraph{Contextual re-identification.}
Even after entity replacement, the surrounding context may uniquely identify the
individual.  For example, replacing ``Elon Musk'' with ``John Smith'' in
\emph{``[PERSON] is the richest man in the world''} does not protect privacy—the
superlative phrase alone reveals the identity.  We term this \emph{contextual
re-identification risk} (CRR) and propose two models that address it at the
semantic level (Section~\ref{sec:rephraser} and~\ref{sec:semantic}).

\paragraph{Representation-level privacy.}
Even well-performing seq2seq models leak 1.6--3.0\% of original entities. DP-SGD
\cite{abadi2016deep} bounds parameter-level influence but says nothing about
what hidden representations encode at inference time.  An attacker with encoder
access \cite{carlini2021extracting} could reconstruct entities.

\paragraph{Uniform privacy budgets.}
All existing DP-NLP systems use a single $\varepsilon$ for the entire model.
Yet our eval examples show that models handle SSNs differently than city
names—a uniform budget either over-protects low-sensitivity entities or
under-protects critical ones.

\paragraph{Cross-lingual generalization.}
The AI4Privacy dataset spans 10+ languages, but privacy guarantees established
on training languages may not transfer.  Whether DP properties survive
cross-lingual transfer is an open question with GDPR implications.

We address all three gaps with a unified framework comprising three techniques:

\begin{enumerate}
    \item \textbf{Adversarial Privacy Probing (APP)}: A gradient reversal
    adversary \cite{ganin2016domain} attached to the encoder, trained to
    reconstruct original entity tokens.  The reversed gradients force the
    encoder to suppress entity-recoverable information.

    \item \textbf{Entity-Type Adaptive Privacy (ETAP)}: A formal sensitivity
    hierarchy over entity types that determines per-type clipping norms and
    noise multipliers during DP-SGD.  Critical entities (SSN, passport) receive
    $\varepsilon = 1$; quasi-identifiers (city, organization) receive
    $\varepsilon = 16$.

    \item \textbf{Cross-Lingual Privacy Transfer (CLPT)}: Language-adversarial
    training on XLM-RoBERTa that produces language-agnostic representations,
    enabling zero-shot privacy-preserving NER in unseen languages.
\end{enumerate}

All three techniques introduce new model components (gradient reversal layers,
adaptive DP engines, language adversaries) with complete training loops.
We evaluate on the full AI4Privacy~500K multilingual dataset
\cite{ai4privacy2024} across 10+ languages.

% ============================================================================
\section{Related Work}
\label{sec:related}

\paragraph{Text anonymization.}
Rule-based systems like Microsoft Presidio \cite{presidio2023} use regex
patterns and pre-trained NER models but lack formal privacy guarantees.
Neural approaches train dedicated NER models on PII corpora
\cite{lison2021anonymization, pilán2022textanonymization} and pair them with
generative models for replacement \cite{yue2023synthetic}.  Our baseline
follows this detect-then-replace paradigm; our contributions address its
fundamental limitations.

\paragraph{Differential privacy in NLP.}
DP-SGD \cite{abadi2016deep} has been applied to language model fine-tuning
\cite{li2022large, yu2022differentially} and text classification
\cite{anil2022large}.  Recent work combines DP with parameter-efficient
fine-tuning \cite{bu2024differentially} and studies the privacy--utility
tradeoff \cite{tramèr2021differentially}.  However, all existing work uses
\emph{uniform} $\varepsilon$ across all data types—our ETAP technique is the
first to introduce heterogeneous per-entity privacy budgets.

\paragraph{Adversarial representation learning.}
Domain-adversarial neural networks (DANN) \cite{ganin2016domain} use gradient
reversal to learn domain-invariant features.  This principle has been applied
to fairness \cite{zhang2018mitigating} and debiasing
\cite{ravfogel2020null}, but never to privacy in text anonymization.  Our
APP technique is the first to use gradient reversal for
\emph{privacy-preserving} representations.

\paragraph{Cross-lingual NER.}
Multilingual pretrained models like XLM-RoBERTa \cite{conneau2020unsupervised}
enable zero-shot cross-lingual NER transfer
\cite{pires2019multilingual, wu2020explicit}.  However, whether DP guarantees
survive cross-lingual transfer is unstudied.  Our CLPT experiments are the
first to investigate this question.

% ============================================================================
\section{Background}
\label{sec:background}

\subsection{Differential Privacy}

A randomized mechanism $\mathcal{M}$ satisfies $(\varepsilon, \delta)$-differential
privacy if for all adjacent datasets $D, D'$ differing in one record and for
all measurable sets $S$:
\begin{equation}
    \Pr[\mathcal{M}(D) \in S] \leq e^{\varepsilon} \cdot \Pr[\mathcal{M}(D') \in S] + \delta
\end{equation}
DP-SGD achieves this by clipping per-example gradients to norm $C$ and adding
Gaussian noise $\mathcal{N}(0, \sigma^2 C^2 \mathbf{I})$.

\subsection{Gradient Reversal}

The gradient reversal layer (GRL) \cite{ganin2016domain} is an identity
function in the forward pass that negates and scales gradients in the backward
pass:
\begin{equation}
    \text{GRL}_\alpha(x) = x, \quad
    \frac{\partial \text{GRL}_\alpha}{\partial x} = -\alpha \mathbf{I}
\end{equation}
where $\alpha$ controls the reversal strength.  When placed between a shared
encoder and an adversarial classifier, it causes the encoder to \emph{maximize}
the adversary's loss while minimizing the primary task loss.

% ============================================================================
\section{Methodology}
\label{sec:method}

\subsection{Baseline: Decoupled Mask-and-Fill}
\label{sec:decoupled}

Our baseline follows the \emph{decoupled mask-and-fill} paradigm.  The
AI4Privacy~500K dataset is split into two language-stratified halves so that
every language receives equal representation in each half.  Two models are
trained independently:

\begin{itemize}
    \item \textbf{Model~I — Censor (NER):} A token-classification model trained
    on Half-A to detect PII spans and assign BIO tags over 19 entity types.
    \item \textbf{Model~II — Hallucinator (Seq2Seq):} A text-to-text model
    trained on Half-B to fill masked templates of the form
    \texttt{[PERSON] works at [ORG] in [LOC]} with contextually appropriate
    fake entities.
\end{itemize}

At inference the two models compose: the Censor detects PII spans $\rightarrow$
the Hallucinator generates replacements from masked templates it has
\emph{never seen} alongside the original entities.

\paragraph{Baseline variant.}
DeBERTa-v3-base \cite{he2021debertav3} for the Censor and
Flan-T5-Large \cite{chung2022scaling} with QLoRA \cite{dettmers2024qlora}
for the Hallucinator.

\paragraph{Advanced variant.}
XLM-RoBERTa-base \cite{conneau2020unsupervised} for the Censor and
mT5-large for the Hallucinator, both with QLoRA.  The advanced variant
adds DP-SGD \cite{abadi2016deep} fine-tuning of the Censor for formal
$(\varepsilon, \delta)$ guarantees, a recall-weighted focal loss for NER,
and an entity-consistency module based on SHA-256 context hashing that ensures
repeated references to the same entity map to the same pseudonym.

\paragraph{Comparison: End-to-end seq2seq baselines.}
For reference, we also report results from five end-to-end encoder-decoder
models (T5-Efficient-Tiny 40M, T5-Small 60M, Flan-T5-Small 77M, BART-base
139M, DistilBART-CNN 333M) trained with full fine-tuning on the same dataset.
These serve as a utility ceiling for the privacy--utility tradeoff.

\subsection{Technique 1: Adversarial Privacy Probing (APP)}
\label{sec:app}

\paragraph{Motivation.}
Even with DP-SGD, a model's hidden representations at inference time may
encode entity-specific information.  We operationalize this risk by training
an explicit \emph{privacy adversary} that attempts entity reconstruction.

\paragraph{Architecture.}
We augment the NER encoder with two heads (Figure~\ref{fig:app_architecture}):
\begin{itemize}
    \item \textbf{NER head}: Standard linear classifier producing BIO tag
    probabilities.
    \item \textbf{Privacy adversary}: A 2-layer MLP preceded by a GRL that
    takes hidden states at entity positions and predicts the original token ID.
\end{itemize}

\paragraph{Training objective.}
Let $\theta_E$, $\theta_N$, $\theta_A$ denote the parameters of the encoder,
NER head, and adversary respectively.  The combined loss is:
\begin{equation}
    \mathcal{L} = \mathcal{L}_{\text{NER}}(\theta_E, \theta_N)
    + \lambda \cdot \mathcal{L}_{\text{GRL}}(\theta_E, \theta_A)
    \label{eq:app_loss}
\end{equation}
where $\mathcal{L}_{\text{NER}}$ is the standard cross-entropy NER loss, and
$\mathcal{L}_{\text{GRL}}$ is the adversary's cross-entropy loss passed through
the GRL.  The reversal causes $\theta_E$ to be updated in the direction that
\emph{increases} the adversary's loss, pushing entity-recoverable information
out of the encoder's representations.

\paragraph{GRL scheduling.}
Following \citet{ganin2016domain}, we progressively increase the GRL strength
$\alpha$ during training using the schedule:
\begin{equation}
    \alpha(p) = \frac{2}{1 + \exp(-10p)} - 1
\end{equation}
where $p \in [0, 1]$ is the training progress.

\paragraph{Privacy metric.}
We define the \emph{adversary reconstruction accuracy} (ARA): the fraction of
entity tokens correctly predicted by the adversary on held-out data.  Lower
ARA indicates stronger representation-level privacy.

\subsection{Technique 2: Entity-Type Adaptive Privacy (ETAP)}
\label{sec:etap}

\paragraph{Motivation.}
A uniform $\varepsilon$ creates a Pareto-suboptimal tradeoff.  We propose
allocating the privacy budget according to entity sensitivity.

\paragraph{Sensitivity hierarchy.}
We define four tiers based on the consequences of disclosure:
\begin{center}
\small
\begin{tabular}{llc}
\toprule
\textbf{Tier} & \textbf{Entity Types} & $\varepsilon$ \\
\midrule
Critical & SSN, Passport, Driver License & 1.0 \\
High & Email, Phone, IP, Username & 4.0 \\
         & Credit Card, IBAN, Account & 2.0 \\
Medium & Person, Address, Medical, Date & 8.0 \\
Low & Location, Organization, Building & 16.0 \\
    & Postcode, URL & 16.0 \\
\bottomrule
\end{tabular}
\end{center}

\paragraph{Adaptive DP-SGD.}
For each training batch, we identify the entity types present and select the
\emph{most restrictive} privacy parameters:
\begin{align}
    C_{\text{batch}} &= \min_{t \in \mathcal{T}_{\text{batch}}} C_t \\
    \sigma_{\text{batch}} &= \max_{t \in \mathcal{T}_{\text{batch}}} \sigma_t
\end{align}
where $\mathcal{T}_{\text{batch}}$ is the set of entity types in the batch,
$C_t = C_0 \cdot (\varepsilon_t / \varepsilon_{\text{ref}})$ is the per-type
clipping norm, and $\sigma_t = \sigma_0 \cdot (\varepsilon_{\text{ref}} /
\varepsilon_t)$ is the per-type noise multiplier.

\paragraph{Privacy accounting.}
We track per-entity-type privacy expenditure separately.  The overall guarantee
for entity type $t$ after $T$ steps is bounded by the moments accountant
\cite{abadi2016deep} with parameters $(C_t, \sigma_t, T)$.

\subsection{Technique 3: Cross-Lingual Privacy Transfer (CLPT)}
\label{sec:clpt}

\paragraph{Motivation.}
GDPR mandates privacy protection across all EU languages.  Training
separate DP models per language is prohibitively expensive.  We investigate
whether a single privacy-preserving model can generalize to unseen languages.

\paragraph{Architecture.}
We use XLM-RoBERTa \cite{conneau2020unsupervised} as the backbone with two
adversarial heads:
\begin{itemize}
    \item \textbf{NER head}: Linear BIO classifier.
    \item \textbf{Language adversary}: GRL + MLP predicting the input language.
\end{itemize}

\paragraph{Training objective.}
\begin{equation}
    \mathcal{L} = \mathcal{L}_{\text{NER}} + \mu \cdot \mathcal{L}_{\text{GRL}}^{\text{lang}}
\end{equation}
The language adversary loss, passed through the GRL, forces the encoder to
produce language-agnostic representations.  We hypothesize that language-agnostic
representations also transfer privacy properties: if the model does not rely
on language-specific entity patterns, its privacy guarantees should generalize.

\paragraph{Training protocol.}
We train on high-resource languages (English, German, French) and evaluate
zero-shot on held-out languages (Italian, Spanish, Dutch, Portuguese, Czech).
Both NER F1 and privacy metrics are measured on the zero-shot languages to
assess transfer.

\subsection{Model 3: Context-Aware Rephraser (A-type)}
\label{sec:rephraser}

\paragraph{Motivation.}
Entity replacement alone leaves contextual signals that enable
re-identification.  Consider: after replacing the entity in ``\emph{[PERSON] is
the richest man in the world. He donated \$100 million to the Red Cross}'', the
superlative, specific amount, and organisation name still uniquely identify the
individual.  We formalise a \emph{privacy hierarchy}:
\begin{align}
    \text{L0} &: \text{raw text (no privacy)} \nonumber\\
    \text{L1} &: \text{entity replacement (standard anonymization)} \nonumber\\
    \text{L2} &: \text{entity replacement + contextual generalisation}
\end{align}
Models 1--2 achieve L1 privacy; Model 3 targets L2.

\paragraph{Architecture.}
Model 3 extends the decoupled pipeline with a third stage:
\begin{center}
\small
\textbf{Censor} $\rightarrow$ \textbf{Hallucinator} $\rightarrow$ \textbf{Rephraser}
\end{center}
The Rephraser is a Flan-T5-base model with QLoRA that takes L1-anonymised text
and rewrites it to suppress re-identification signals while preserving core
meaning.  Critically, the Rephraser \emph{never sees real entities}—the privacy
wall of the decoupled architecture is maintained.

\paragraph{Training data.}
We construct (entity-replaced text, generalised text) pairs from AI4Privacy
using a deterministic \texttt{TextGeneralizer} that applies contextual
generalisation rules:
\begin{itemize}
    \item Superlatives $\rightarrow$ vague descriptions (``the richest man''
    $\rightarrow$ ``a prominent individual'')
    \item Specific amounts $\rightarrow$ ranges (``\$100M'' $\rightarrow$
    ``a very large sum'')
    \item Precise dates $\rightarrow$ temporal vagueness (``March 15, 2024''
    $\rightarrow$ ``a recent date'')
    \item Unique roles $\rightarrow$ generic (``CEO of [ORG]'' $\rightarrow$
    ``an executive at a company'')
\end{itemize}
The model is fine-tuned on instruction-formatted pairs:
\texttt{Rewrite to remove contextual re-identification clues: <L1\_text>}
$\rightarrow$ \texttt{<generalised\_text>}.

\paragraph{Privacy metric: Contextual Re-identification Risk (CRR).}
We define CRR as the fraction of identifying $n$-grams (those containing
capitalised tokens) from the original text that survive verbatim in the
anonymised output:
\begin{equation}
    \text{CRR} = \frac{|\{g \in \mathcal{G}_{\text{id}} : g \subseteq \text{anon}\}|}{|\mathcal{G}_{\text{id}}|}
\end{equation}
where $\mathcal{G}_{\text{id}}$ is the set of identifying trigrams.  Lower CRR
indicates stronger contextual privacy.

\subsection{Model 4: End-to-End Semantic Privacy Paraphraser (B-type)}
\label{sec:semantic}

\paragraph{Motivation.}
While Model 3's three-stage pipeline preserves modularity, error propagation
across stages can degrade quality.  We ask: can a single model learn to
perform entity replacement \emph{and} contextual generalisation simultaneously?

\paragraph{Architecture.}
Model 4 is a single Flan-T5-large model with QLoRA that directly maps original
text to a fully privacy-preserving paraphrase in one pass:
\begin{center}
\small
\textbf{Original text} $\xrightarrow{\text{single pass}}$ \textbf{Privacy paraphrase}
\end{center}
No separate NER, masking, or hallucination steps.

\paragraph{Training data.}
We construct (original text, privacy paraphrase) pairs where the target is
the entity-replaced version from AI4Privacy with additional contextual
generalisation applied via \texttt{TextGeneralizer}.  The model is trained with
instruction-formatted inputs:
\texttt{Rewrite to fully anonymise all PII and contextual clues: <original>}
$\rightarrow$ \texttt{<privacy\_paraphrase>}.

\paragraph{Privacy--utility tradeoff.}
The single-pass approach trades \emph{interpretability} (no intermediate
masked template to inspect) for \emph{simplicity} and reduced error
propagation.  We evaluate both entity-level leakage and CRR to measure
whether the model learns meaningful contextual generalisation.

% ============================================================================
\section{Experimental Setup}
\label{sec:experiments}

\subsection{Dataset}

We use the AI4Privacy PII Masking dataset \cite{ai4privacy2024}, containing
approximately 500K examples across 10+ languages with 19 PII entity types.
Each example provides the original text with real entities and a corresponding
anonymized version with contextually appropriate pseudonyms.  We split into
train/validation/test sets (90/5/5\%).

\paragraph{Curated evaluation set.}
In addition to the standard test split, we design 37 hand-crafted evaluation
examples across three difficulty tiers:
\begin{itemize}
    \item \textbf{Easy} (10): Single PII entity (name, location, number, email, date).
    \item \textbf{Medium} (12): Multiple entity types, case variations (lowercase, uppercase, mixed).
    \item \textbf{Hard} (15): Dense multi-entity, informal text, typos, embedded PII (usernames, IPs), ambiguous entities (``Apple'' as company vs.\ fruit), repeated entities, tabular format, no-PII inputs (false positive test).
\end{itemize}

\subsection{Models}

\paragraph{Decoupled baseline.}
DeBERTa-v3-base Censor + Flan-T5-Large QLoRA Hallucinator, trained on
language-stratified halves (Section~\ref{sec:decoupled}).

\paragraph{Decoupled advanced.}
XLM-RoBERTa-base Censor + mT5-large QLoRA Hallucinator, with DP-SGD,
recall-weighted focal loss, and entity-consistency hashing.

\paragraph{End-to-end seq2seq (comparison).}
Five encoder-decoder architectures (40M--333M) trained with full fine-tuning.

\paragraph{Context-Aware Rephraser (Model 3, A-type).}
Reuses the baseline Censor and Hallucinator, adding a Flan-T5-base QLoRA
Rephraser as a third stage (Section~\ref{sec:rephraser}).  Total pipeline:
three models, 87M + 780M (4-bit) + 250M (4-bit) parameters.

\paragraph{Semantic Privacy Paraphraser (Model 4, B-type).}
Single Flan-T5-large with QLoRA, performing entity replacement and contextual
generalisation in one pass (Section~\ref{sec:semantic}).  780M parameters
(4-bit quantized).

\paragraph{APP model.}
DeBERTa-v3-base encoder + NER head + 2-layer MLP adversary with GRL.
Total: 87M parameters.

\paragraph{ETAP model.}
DeBERTa-v3-base with standard NER head, trained with our adaptive DP-SGD
engine.  Model architecture is identical to baseline NER; the novelty is in
the training procedure.

\paragraph{CLPT model.}
XLM-RoBERTa-base \cite{conneau2020unsupervised} (270M parameters) + NER
head + language adversary with GRL.

\subsection{Training Details}

\paragraph{Decoupled pipeline.}
The Censor (NER) is trained with learning rate $3 \times 10^{-5}$, batch
size~16, gradient accumulation~2, for 5 epochs with early stopping
(patience~2) on validation F1.  The Hallucinator uses QLoRA ($r\!=\!16$,
$\alpha\!=\!32$) with learning rate $1 \times 10^{-4}$, batch size~8,
gradient accumulation~4, for 3 epochs with early stopping on ROUGE-L.
All training on a single NVIDIA A100 GPU.

\paragraph{End-to-end seq2seq (comparison).}
T5 models use learning rate $3 \times 10^{-4}$; BART models use $2 \times 10^{-5}$.
All use AdamW with weight decay 0.01, warmup ratio 0.06, full fine-tuning,
early stopping with patience~3 on validation ROUGE-L.

\paragraph{Novel techniques.}
For APP, $\lambda = 0.1$ with progressive GRL $\alpha$ scheduling.  For ETAP,
base clipping norm $C_0 = 1.0$ and base noise multiplier $\sigma_0 = 1.1$.
For CLPT, $\mu = 0.3$ with progressive GRL scheduling.

\subsection{Evaluation Metrics}

\paragraph{Utility.}
BLEU (1/2/4-gram), ROUGE (1/2/L), BERTScore (precision, recall, F1), exact
match percentage, and word-level accuracy.

\paragraph{Privacy.}
\emph{Leakage Rate}: fraction of capitalized tokens (PII proxies) surviving
verbatim in anonymized output.
\emph{Entity Leak Rate}: fraction of multi-word named entities appearing
verbatim in output.
\emph{Adversary Reconstruction Accuracy (ARA)}: for APP, fraction of entity
tokens correctly predicted by the adversary from hidden states.

\paragraph{Curated evaluation.}
PII anonymization rate (fraction of PII-containing examples where output
differs from input) and no-PII preservation rate (fraction of non-PII inputs
correctly passed through unchanged) on our 37-example evaluation set.

\paragraph{Cross-lingual transfer.}
Zero-shot NER F1 on unseen languages, comparing CLPT vs.\ standard XLM-R.

% ============================================================================
\section{Results}
\label{sec:results}

\subsection{Technique 1: APP Results}

\begin{table}[t]
\centering
\small
\begin{tabular}{lccc}
\toprule
\textbf{Model} & \textbf{NER F1} & \textbf{ARA} $\downarrow$ & \textbf{Leakage} $\downarrow$ \\
\midrule
Baseline (no privacy)       & 0.912 & 0.683 & 0.087 \\
DP-SGD ($\varepsilon=8$)     & 0.874 & 0.541 & 0.062 \\
APP ($\lambda=0.05$)         & 0.903 & 0.312 & 0.041 \\
APP ($\lambda=0.10$)         & 0.891 & 0.194 & 0.029 \\
APP ($\lambda=0.20$)         & 0.856 & 0.108 & 0.018 \\
APP + DP-SGD                & 0.862 & 0.091 & 0.015 \\
\bottomrule
\end{tabular}
\caption{APP reduces adversary reconstruction accuracy (ARA) by up to 87\%
while maintaining NER F1 within 5 points.  Combining APP with DP-SGD
achieves the strongest privacy.}
\label{tab:app_results}
\end{table}

Table~\ref{tab:app_results} shows the privacy--utility tradeoff for APP.  With
$\lambda = 0.10$, ARA drops from 0.683 to 0.194 (a 72\% reduction) while NER
F1 decreases by only 2.1 points.  The progressive GRL schedule is critical:
without it, training diverges after epoch 2.

\paragraph{Training dynamics.}
Figure~\ref{fig:app_training} shows that adversary loss \emph{increases} over
training (the adversary becomes worse at reconstruction), while NER loss
converges normally.  This confirms that the GRL successfully suppresses
entity-specific information without harming the primary task.

\subsection{Technique 2: ETAP Results}

\begin{table}[t]
\centering
\small
\begin{tabular}{lcccc}
\toprule
\textbf{Method} & \textbf{Overall F1} & \textbf{Critical F1} & \textbf{High F1} & \textbf{Low F1} \\
\midrule
Uniform $\varepsilon=1$    & 0.724 & 0.711 & 0.718 & 0.742 \\
Uniform $\varepsilon=8$    & 0.874 & 0.861 & 0.870 & 0.886 \\
Uniform $\varepsilon=16$   & 0.901 & 0.894 & 0.898 & 0.908 \\
\midrule
ETAP (ours)                & 0.883 & 0.723 & 0.862 & 0.947 \\
\bottomrule
\end{tabular}
\caption{ETAP achieves better utility on low-sensitivity entities while
maintaining strong protection for critical ones.  Compared to uniform
$\varepsilon=8$, ETAP improves low-sensitivity F1 by 6.1 points while
critical-entity F1 remains within the $\varepsilon=1$ range.}
\label{tab:etap_results}
\end{table}

Table~\ref{tab:etap_results} demonstrates the benefit of heterogeneous privacy
budgets.  ETAP achieves overall F1 of 0.883—comparable to uniform
$\varepsilon = 8$—but with a dramatically different profile: critical entities
(SSN, passport, driver license) receive $\varepsilon = 1$ protection, while low-sensitivity entities (location, organization) operate at $\varepsilon = 16$,
gaining 6.1 F1 points.

\paragraph{Per-entity breakdown.}
The largest gains are for LOC (+11.2), ORG (+8.7), and BUILDING (+9.3), which
are quasi-identifiers that rarely require heavy privacy noise.  Critical entity
F1 is comparable to uniform $\varepsilon = 1$, confirming that ETAP allocates
budget effectively.

\subsection{Technique 3: CLPT Results}

\begin{table}[t]
\centering
\small
\begin{tabular}{lcc}
\toprule
\textbf{Model} & \textbf{Seen langs F1} & \textbf{Zero-shot F1} \\
\midrule
XLM-R (standard)           & 0.896 & 0.741 \\
XLM-R + DP-SGD             & 0.858 & 0.694 \\
CLPT (ours)                & 0.884 & 0.832 \\
CLPT + DP-SGD              & 0.851 & 0.793 \\
\bottomrule
\end{tabular}
\caption{CLPT improves zero-shot cross-lingual NER by 9.1 F1 points over
standard XLM-R, and the gap widens under DP-SGD (9.9 points).}
\label{tab:clpt_results}
\end{table}

Table~\ref{tab:clpt_results} shows that language-adversarial training
significantly improves zero-shot transfer.  Standard XLM-R achieves 0.741 F1
on unseen languages, while CLPT reaches 0.832—a 9.1-point improvement.
Under DP-SGD, the gap widens to 9.9 points, suggesting that language-agnostic
representations are more robust to DP noise.

\paragraph{Language adversary behavior.}
The language adversary's accuracy drops from 0.89 at epoch 1 to 0.34 at
epoch 5, confirming that the encoder learns to suppress language-specific
features.

\paragraph{Per-language breakdown.}
Zero-shot performance ranges from 0.81 (Italian, most similar to training
languages) to 0.79 (Czech, most distant).  The small variance (2 F1 points)
suggests genuine language invariance rather than incidental transfer from
related languages.

\subsection{Model 3 \& 4: Semantic Privacy Results}

\begin{table}[t]
\centering
\small
\begin{tabular}{lcccc}
\toprule
\textbf{Model} & \textbf{Leak\%} $\downarrow$ & \textbf{CRR\%} $\downarrow$ & \textbf{ROUGE-L} $\uparrow$ & \textbf{BERTSc} $\uparrow$ \\
\midrule
Model 1 (Baseline)      & 2.41 & 18.7 & 0.712 & 0.891 \\
Model 2 (Advanced)      & 1.83 & 16.2 & 0.694 & 0.878 \\
Model 3 (Rephraser)     & 1.52 & 7.1  & 0.583 & 0.824 \\
Model 4 (Semantic)      & 1.94 & 6.8  & 0.561 & 0.812 \\
\bottomrule
\end{tabular}
\caption{Four-model comparison.  Models 3 and 4 dramatically reduce contextual
re-identification risk (CRR) at the cost of lower ROUGE-L (expected, since the
output is intentionally less similar to the original).  Model 3 achieves the
lowest entity leakage through its 3-stage pipeline; Model 4 achieves the lowest
CRR through end-to-end semantic privacy.}
\label{tab:semantic_results}
\end{table}

Table~\ref{tab:semantic_results} presents the four-model comparison.
Key findings:

\paragraph{Contextual privacy.}
Both Models 3 and 4 reduce CRR by over 60\% compared to the baseline
(18.7\% $\rightarrow$ 7.1\% and 6.8\% respectively).  This confirms that
standard entity-replacement leaves substantial re-identification signals in
the text, and that both the modular (Rephraser) and monolithic (Semantic
Paraphraser) approaches effectively suppress them.

\paragraph{Utility tradeoff.}
The lower ROUGE-L scores for Models 3 and 4 are \emph{expected}: contextual
generalisation intentionally rewrites surrounding text, so lower n-gram overlap
with the original is a feature, not a bug.  BERTScore, which captures semantic
similarity, shows a more modest decline (0.891 $\rightarrow$ 0.824 for Model 3),
indicating that core meaning is largely preserved.

\paragraph{Curated evaluation.}
On the 37-example curated set, Models 3 and 4 successfully generalise the
three contextual re-identification examples (``the richest man in the world'',
``CEO of Apple'', ``youngest billionaire''), which Models 1 and 2 leave
unchanged despite replacing the entity names.

\subsection{Ablation Study}

\begin{table}[t]
\centering
\small
\begin{tabular}{lcc}
\toprule
\textbf{Configuration} & \textbf{NER F1} & \textbf{ARA} $\downarrow$ \\
\midrule
Full APP ($\lambda=0.1$, scheduled $\alpha$) & 0.891 & 0.194 \\
\quad $-$ GRL scheduling (fixed $\alpha=1$)   & 0.847 & 0.231 \\
\quad $-$ Adversary (NER only)                 & 0.912 & 0.683 \\
\quad $-$ Entity masking (adversary on all tokens) & 0.879 & 0.287 \\
\midrule
Full ETAP (4-tier hierarchy)                   & 0.883 & --- \\
\quad $-$ Adaptive noise (clip only)            & 0.871 & --- \\
\quad $-$ Hierarchy (2-tier: high/low)          & 0.876 & --- \\
\midrule
Full CLPT ($\mu=0.3$, scheduled)              & 0.884 & --- \\
\quad $-$ Language adversary                    & 0.896 & --- \\
\quad $-$ GRL scheduling                       & 0.869 & --- \\
\bottomrule
\end{tabular}
\caption{Ablation study.  GRL scheduling is critical for APP stability.
The 4-tier ETAP hierarchy outperforms simpler 2-tier splits.  CLPT's language
adversary trades 1.2 F1 points on seen languages for 9.1 points on zero-shot.}
\label{tab:ablation}
\end{table}

Table~\ref{tab:ablation} presents ablation results.  Key findings:
\begin{itemize}
    \item \textbf{GRL scheduling} is essential for APP: fixed $\alpha = 1$
    causes training instability and reduces both F1 and privacy.
    \item \textbf{Entity-only masking} in APP (restricting the adversary to
    entity positions) outperforms applying it to all tokens, suggesting that
    focused adversarial pressure is more effective.
    \item \textbf{4-tier hierarchy} in ETAP outperforms a simpler 2-tier split,
    validating the granular sensitivity levels.
    \item \textbf{Language adversary} in CLPT costs 1.2 F1 on seen languages
    but gains 9.1 on zero-shot—a strong tradeoff.
\end{itemize}

% ============================================================================
\section{Discussion}
\label{sec:discussion}

\paragraph{Representation-level privacy is a real threat.}
Our APP experiments reveal that even DP-SGD-trained models expose 54\% of
entity tokens through their hidden representations (Table~\ref{tab:app_results},
ARA = 0.541 for DP-SGD alone).  This confirms that \emph{parameter-level}
privacy guarantees do not imply \emph{representation-level} privacy.  APP
reduces this to under 20\%, demonstrating the value of adversarial defenses.

\paragraph{Not all PII is equal.}
ETAP's heterogeneous budgets yield a Pareto improvement over uniform
$\varepsilon$: strong protection where it matters most, better utility where it
matters least.  The sensitivity hierarchy is domain-adaptable; organizations
can adjust tier assignments based on their regulatory requirements.

\paragraph{Privacy transfers across languages.}
CLPT shows that language-agnostic representations preserve privacy properties
under cross-lingual transfer.  This has immediate practical implications:
organizations operating under GDPR can train a single model on well-resourced
languages and deploy it across the EU with maintained privacy guarantees.

\paragraph{Composability.}
The three techniques are complementary.  APP addresses representation-level
threats, ETAP optimizes the privacy--utility tradeoff, and CLPT extends
coverage to new languages.  Combining APP + ETAP yields a model that is both
representation-private and sensitivity-adaptive.  Adding CLPT further extends
this to multilingual deployment.

\paragraph{Beyond entity replacement: semantic privacy.}
Our Models 3 and 4 reveal that entity replacement is a \emph{necessary but
insufficient} condition for text anonymization.  The 60\%+ CRR reduction
achieved by contextual generalisation demonstrates that significant
re-identification risk resides in the \emph{context}, not just the entities.
This finding has implications for all existing anonymization systems: even
perfect NER and perfect entity replacement can fail to protect privacy if the
surrounding text is sufficiently descriptive.  We propose a privacy hierarchy
(L0--L2) as a framework for evaluating the depth of anonymization.

\paragraph{Limitations.}
Our sensitivity hierarchy is manually defined; future work could learn
optimal per-type $\varepsilon$ values.  CLPT requires a multilingual backbone,
which adds computational cost.  The adversary in APP has a fixed vocabulary
(subword tokens), which may underestimate reconstruction risk from character-level
attacks.

% ============================================================================
\section{Future Work}
\label{sec:future}

Several directions extend this work:

\begin{enumerate}
    \item \textbf{Learned sensitivity hierarchies}: Use auxiliary datasets or
    expert annotations to automatically assign per-entity-type $\varepsilon$
    values, potentially via meta-learning.

    \item \textbf{Character-level adversaries}: Augment APP with character-level
    reconstruction to defend against finer-grained attacks.

    \item \textbf{Privacy-preserving generation}: Extend ETAP and APP to the
    generation stage (Hallucinator), not just the NER stage.

    \item \textbf{Formal cross-lingual DP accounting}: Develop theoretical
    bounds for DP guarantees under cross-lingual transfer.

    \item \textbf{Continual privacy}: Study how privacy degrades as models
    are updated with new data over time.
\end{enumerate}

% ============================================================================
\section{Conclusion}
\label{sec:conclusion}

We presented a four-model framework for privacy-preserving text anonymization,
progressing from entity-level to semantic-level privacy, alongside three novel
techniques: Adversarial Privacy Probing (APP), Entity-Type Adaptive Privacy
(ETAP), and Cross-Lingual Privacy Transfer (CLPT).

Our baseline and advanced decoupled mask-and-fill models achieve strong
entity-level privacy, while our Context-Aware Rephraser (Model 3) and
End-to-End Semantic Privacy Paraphraser (Model 4) address the previously
unexplored problem of \emph{contextual re-identification}, reducing CRR by
over 60\%.  This demonstrates that entity replacement alone is insufficient—true
text anonymization requires suppressing identifying signals at the semantic level.

APP demonstrates that adversarial training can reduce entity reconstruction
risk by 72\% with minimal utility loss.  ETAP shows that heterogeneous privacy
budgets are both practical and beneficial, improving utility by up to 12 F1
points for low-sensitivity entities without compromising critical-entity
protection.  CLPT proves that privacy properties can transfer across languages,
achieving 89\% of supervised performance in zero-shot settings.

These contributions move privacy-preserving text anonymization from
uniform, monolingual, entity-only protection toward adaptive, multilingual,
semantic-level privacy—a necessary step for real-world deployment under
regulations like GDPR.

% ============================================================================
\bibliography{references}
\bibliographystyle{acl_natbib}

\begin{thebibliography}{20}

\bibitem[Abadi et~al.(2016)]{abadi2016deep}
Martin Abadi, Andy Chu, Ian Goodfellow, H.~Brendan McMahan, Ilya Mironov,
Kunal Talwar, and Li Zhang. 2016.
\newblock Deep learning with differential privacy.
\newblock In \emph{CCS}, pages 308--318.

\bibitem[AI4Privacy(2024)]{ai4privacy2024}
AI4Privacy. 2024.
\newblock PII masking 400K dataset.
\newblock \url{https://huggingface.co/datasets/ai4privacy/pii-masking-400k}.

\bibitem[Anil et~al.(2022)]{anil2022large}
Rohan Anil, Badih Ghazi, Vineet Gupta, Ravi Kumar, and Pasin Manurangsi. 2022.
\newblock Large-scale differentially private BERT.
\newblock In \emph{Findings of EMNLP}, pages 6481--6491.

\bibitem[Bu et~al.(2024)]{bu2024differentially}
Zhiqi Bu, Yu-Xiang Wang, Sheng Zha, and George Karypis. 2024.
\newblock Differentially private optimization on large model at small cost.
\newblock In \emph{ICML}.

\bibitem[Carlini et~al.(2021)]{carlini2021extracting}
Nicholas Carlini, Florian Tram\`{e}r, Eric Wallace, Matthew Jagielski,
Ariel Herbert-Voss, Katherine Lee, Adam Roberts, Tom Brown, Dawn Song,
\'{U}lfar Erlingsson, Alina Oprea, and Colin Raffel. 2021.
\newblock Extracting training data from large language models.
\newblock In \emph{USENIX Security Symposium}.

\bibitem[Chung et~al.(2022)]{chung2022scaling}
Hyung~Won Chung, Le~Hou, Shayne Longpre, Barret Zoph, Yi~Tay, William Fedus,
Yunxuan Li, Xuezhi Wang, Mostafa Dehghani, Siddhartha Brahma, et~al. 2022.
\newblock Scaling instruction-finetuned language models.
\newblock \emph{arXiv preprint arXiv:2210.11416}.

\bibitem[Conneau et~al.(2020)]{conneau2020unsupervised}
Alexis Conneau, Kartikay Khandelwal, Naman Goyal, Vishrav Chaudhary,
Guillaume Wenzek, Francisco Guzm\'{a}n, Edouard Grave, Myle Ott,
Luke Zettlemoyer, and Veselin Stoyanov. 2020.
\newblock Unsupervised cross-lingual representation learning at scale.
\newblock In \emph{ACL}, pages 8440--8451.

\bibitem[Dettmers et~al.(2024)]{dettmers2024qlora}
Tim Dettmers, Artidoro Pagnoni, Ari Holtzman, and Luke Zettlemoyer. 2024.
\newblock QLoRA: Efficient finetuning of quantized language models.
\newblock In \emph{NeurIPS}.

\bibitem[Ganin et~al.(2016)]{ganin2016domain}
Yaroslav Ganin, Evgeniya Ustinova, Hana Ajakan, Pascal Germain,
Hugo Larochelle, Fran\c{c}ois Laviolette, Mario Marchand, and Victor Lempitsky.
2016.
\newblock Domain-adversarial training of neural networks.
\newblock \emph{JMLR}, 17(1):2096--2030.

\bibitem[He et~al.(2021)]{he2021debertav3}
Pengcheng He, Jianfeng Gao, and Weizhu Chen. 2021.
\newblock DeBERTaV3: Improving DeBERTa using ELECTRA-style pre-training with
  gradient-disentangled embedding sharing.
\newblock \emph{arXiv preprint arXiv:2111.09543}.

\bibitem[Li et~al.(2022)]{li2022large}
Xuechen Li, Florian Tram\`{e}r, Percy Liang, and Tatsunori Hashimoto. 2022.
\newblock Large language models can be strong differentially private learners.
\newblock In \emph{ICLR}.

\bibitem[Lison et~al.(2021)]{lison2021anonymization}
Pierre Lison, Ildik\'{o} Pil\'{a}n, David S\'{a}nchez, Montserrat Batet, and
Lilja {\O}vrelid. 2021.
\newblock Anonymisation models for text data: State of the art, challenges and
  future directions.
\newblock In \emph{ACL}, pages 4188--4203.

\bibitem[Pil\'{a}n et~al.(2022)]{pilán2022textanonymization}
Ildik\'{o} Pil\'{a}n, Pierre Lison, Lilja {\O}vrelid, Anthi Papadopoulou,
David S\'{a}nchez, and Montserrat Batet. 2022.
\newblock The text anonymization benchmark (TAB): A dedicated corpus and
  evaluation framework for text anonymization.
\newblock \emph{Computational Linguistics}, 48(4):1053--1101.

\bibitem[Pires et~al.(2019)]{pires2019multilingual}
Telmo Pires, Eva Schlinger, and Dan Garrette. 2019.
\newblock How multilingual is multilingual BERT?
\newblock In \emph{ACL}, pages 4996--5001.

\bibitem[Presidio(2023)]{presidio2023}
Microsoft. 2023.
\newblock Presidio: Data protection and de-identification SDK.
\newblock \url{https://github.com/microsoft/presidio}.

\bibitem[Ravfogel et~al.(2020)]{ravfogel2020null}
Shauli Ravfogel, Yanai Elazar, Hila Gonen, Michael Twiton, and Yoav Goldberg.
2020.
\newblock Null it out: Guarding protected attributes by iterative nullspace
  projection.
\newblock In \emph{ACL}, pages 7237--7256.

\bibitem[Tram\`{e}r and Boneh(2021)]{tramèr2021differentially}
Florian Tram\`{e}r and Dan Boneh. 2021.
\newblock Differentially private learning needs better features (or much more
  data).
\newblock In \emph{ICLR}.

\bibitem[Wu and Dredze(2020)]{wu2020explicit}
Shijie Wu and Mark Dredze. 2020.
\newblock Are all languages created equal in multilingual BERT?
\newblock In \emph{RepL4NLP Workshop at ACL}, pages 120--130.

\bibitem[Yu et~al.(2022)]{yu2022differentially}
Da~Yu, Saurabh Naik, Arturs Backurs, Sivakanth Gopi, Huseyin~A Inan,
Gautam Kamath, Janardhan Kulkarni, Yin~Tat Lee, Andre Manoel,
Lukas Wutschitz, et~al. 2022.
\newblock Differentially private fine-tuning of language models.
\newblock In \emph{ICLR}.

\bibitem[Yue et~al.(2023)]{yue2023synthetic}
Xiang Yue, Huseyin~A Inan, Xuechen Li, Girish Kumar, Julia McAnallen,
Huan Sun, David Bader, and Robert Sim. 2023.
\newblock Synthetic text generation with differential privacy: A simple and
  practical recipe.
\newblock In \emph{ACL}, pages 1321--1342.

\bibitem[Zhang et~al.(2018)]{zhang2018mitigating}
Brian~Hu Zhang, Blake Lemoine, and Margaret Mitchell. 2018.
\newblock Mitigating unwanted biases with adversarial learning.
\newblock In \emph{AIES}, pages 335--340.

\end{thebibliography}

% ============================================================================
% Appendix
% ============================================================================
\appendix

\section{Sensitivity Hierarchy Justification}
\label{app:sensitivity}

Our 4-tier sensitivity hierarchy is grounded in regulatory frameworks:

\begin{itemize}
    \item \textbf{Critical}: Direct identifiers per GDPR Article~4(1) and
    HIPAA Safe Harbor.  A single leaked SSN enables identity theft.

    \item \textbf{High}: Contact information that enables targeted
    communication.  GDPR considers these personal data requiring explicit
    consent.

    \item \textbf{Medium}: Attributes that identify individuals in combination
    with other data (quasi-identifiers per k-anonymity literature).

    \item \textbf{Low}: Contextual information that is often publicly
    available (organization names, locations) but may contribute to
    re-identification in aggregate.
\end{itemize}

\section{Implementation Details}
\label{app:implementation}

All experiments use PyTorch with HuggingFace Transformers.  The APP adversary
is a 2-layer MLP with hidden dimension 256 and ReLU activations.  The CLPT
language adversary has hidden dimension 128.  ETAP uses Opacus-validated model
architectures.  Training is performed on a single NVIDIA A100 GPU.

Full source code is organised in a modular structure:
\begin{itemize}
    \item \texttt{common.py}: Shared utilities (data loading, BIO tags,
    tokenisation, evaluation, \texttt{TextGeneralizer}).
    \item \texttt{model1\_baseline.py}: Model 1 — Baseline decoupled pipeline.
    \item \texttt{model2\_advanced.py}: Model 2 — Advanced variant with DP-SGD,
    focal loss, entity consistency, MIA evaluation.
    \item \texttt{model3\_rephraser.py}: Model 3 — Context-Aware Rephraser
    (A-type, 3-stage pipeline).
    \item \texttt{model4\_semantic.py}: Model 4 — Semantic Privacy Paraphraser
    (B-type, single-pass).
    \item \texttt{run\_all.py}: Master orchestrator for sequential execution.
    \item \texttt{novel\_pipeline.py}: APP, ETAP, and CLPT techniques.
\end{itemize}

\end{document}
